\documentclass[12pt]{article}
\usepackage[utf8]{inputenc}
\usepackage[T1]{fontenc}
\usepackage{amsmath,amssymb}
\usepackage{geometry}
\geometry{margin=2.5cm}

\begin{document}

\noindent\underline{Dem}:

\bigskip
\noindent (1) Pp.\ $a_i^{-1}\equiv_d a_j^{-1} \pmod H \Rightarrow a_i^{-1}a_j\in H \Rightarrow a_i\equiv_s a_j \pmod H$ \ $\bot$.

\bigskip
\noindent (2) $x\in G \Rightarrow \exists\, i\in I$ a.î.\ $x^{-1}\equiv_s a_i \pmod H \Rightarrow (x^{-1})^{-1}a_i\in H \Rightarrow$

\noindent şi $xa_i\in H \Rightarrow x(a_i^{-1})^{-1}\in H \Rightarrow x\equiv_d a_i^{-1} \pmod H$.

\bigskip
\noindent $G$ grup, \ $f:G\to G$, \ $f(x)=x^{-1}$ = bijecție.

\bigskip
\noindent
\[
G = \bigcup_{i\in I} a_iH
\]

\noindent $|G|<\infty \Rightarrow |I|=d<\infty$

\medskip
\noindent Pp.\ că $a_1,a_2,\ldots,a_d$ este o transversală [la] st.\ a lui $H$ în $G$.

\[
|G|=n \qquad |H|=m
\]

\noindent $|G| = |H|[G:H]$ -- Th.\ lui Lagrange

\noindent $|H|=|aH|$.

\[
|G| = \sum_{i=1}^{d} |a_iH| = \underbrace{|H|+\cdots+|H|}_{d} = d|H| = dm.
\]

\noindent $f : H \to aH$

\noindent $f$ surj.

\noindent $f(x)=ax$; \quad $ax_1=ax_2 \Rightarrow x_1=x_2$ \quad $\Rightarrow$ \ $f$ bij.

\bigskip
\noindent\underline{Exp}: de grup cu 12 elem.\ ce admite un subgrup de 4 elemente.

\[
(\mathbb{Z}_{12},+)
\]
\[
\hat 0,\hat 1,\hat 2,\hat 3,\hat 4,\hat 5,\hat 6,\hat 7,\hat 8,\hat 9,\hat{10},\hat{11}
\]
\[
\{\hat 0,\hat 3,\hat 6,\hat 9\} \leq (\mathbb{Z}_{12},+).
\]

\end{document}
